\section{Detection-safety model}
\label{sec:safety}

The second constraint of the sizing problem concerns the ability to notice that
a subset of validators has stopped behaving correctly.

We claim no novelty for the inequality used here. It is Hoeffding's bound with
a design-effect correction for correlated observations, both textbook material
\cite{hoeffding1963,boucheron2013}. The section exists for two reasons. First,
Theorem~\ref{thm:window} needs a constraint that is monotone in $\n$ in the
opposite direction to throughput, and this bound supplies one. Second, the
quantities entering it --- the detection probability, the false-positive rate
and the correlation between faulty validators --- were \emph{assumed} in our
earlier work \cite{peniak2026b}; here they are estimated from labelled
fault-injection runs. The contribution is the substitution of measurement for
assumption, not the inequality.

\subsection{Fault and detection model}

\begin{Def}[Fault model]
\label{def:faults}
Of $\n$ validators, $\kf$ exhibit \emph{omission and duplication} faults: a
faulty validator drops a fraction of its outgoing consensus traffic and
duplicates another fraction. This is a strict subset of Byzantine behaviour; it
excludes equivocation and forged signatures.
\end{Def}

\begin{Def}[Participation detector]
\label{def:detector}
Over a window of $\w$ consecutive rounds the monitor observes, for each
validator $v_{i}$, the indicator $X_{i}\in\{0,1\}$ that $v_{i}$ failed to meet
its expected participation quota. The estimator of the faulty fraction is
$\Bhat=\n^{-1}\sum_{i=1}^{\n}X_{i}$, with
\begin{equation}
\mu_{B}\eqdef\E[\Bhat]=\pd\frac{\kf}{\n}+\pf\frac{\n-\kf}{\n},
\label{eq:muB}
\end{equation}
where $\pd$ is the detection probability for a faulty validator and $\pf$ the
false-positive rate for a correct one.
\end{Def}

Both $\pd$ and $\pf$ are empirical properties of the monitor and the network,
and are estimated in Section~\ref{sec:results} from ground-truth-labelled
fault-injection runs.

\subsection{Miss probability}

Let $\thr$ be the alarm threshold, and let the effective margin be
$\varepsilon_{0}\eqdef\mu_{B}-\thr>0$. A \emph{miss} is the event
$\Bhat\le\thr$ while the true faulty fraction exceeds it.

\begin{Lemma}[Miss bound under independence]
\label{lem:hoeffding}
If the $X_{i}$ are independent, then over $\w$ windows
$\pmiss=\Prob(\Bhat-\mu_{B}\le-\varepsilon_{0})\le
\exp(-2\n\w\varepsilon_{0}^{2})$.
\end{Lemma}

\begin{Proof}
Hoeffding's inequality \cite{hoeffding1963} applied to the mean of $\n\w$
independent variables bounded in $[0,1]$.
\end{Proof}

Independence is the wrong assumption when faults are coordinated, which is the
case we care about. We therefore replace $\n$ by an effective sample size.

\begin{Assumption}[Exchangeable correlation]
\label{as:rho}
The indicators of the $\kf$ faulty validators are exchangeable with pairwise
correlation $\rho\in[0,1)$; indicators of correct validators are independent.
\end{Assumption}

Assumption~\ref{as:rho} is a modelling convenience. Real coordinated faults need
not be exchangeable, and $\rho$ may drift with the attack intensity. What makes
it tolerable here is that $\rho$ is measured per injection level rather than
fixed a priori, so a violation shows up as instability of the estimate across
levels.

\begin{State}[Miss bound under correlation]
\label{st:rho}
Under Assumption~\ref{as:rho},
\begin{equation}
\pmiss(\n,\w)\le\exp\bigl(-2\,\neff\,\w\,\varepsilon_{0}^{2}\bigr),
\qquad
\neff=\frac{\n}{1+(\kf-1)\rho}.
\label{eq:pmiss}
\end{equation}
\end{State}

\begin{Proof}
The variance of the mean of $\kf$ exchangeable indicators with pairwise
correlation $\rho$ is inflated by the design effect $1+(\kf-1)\rho$ relative to
the independent case. Substituting the variance-matched independent sample size
$\neff=\n/(1+(\kf-1)\rho)$ into the sub-Gaussian tail bound of
Lemma~\ref{lem:hoeffding} gives~\eqref{eq:pmiss}. The substitution is exact for
the variance and an approximation for the tail; it is conservative whenever the
indicator average remains sub-Gaussian with the matched variance proxy.
\end{Proof}

\begin{Corollary}[Monotonicity in $\n$]
\label{cor:pmiss-monotone}
For fixed $\w$, $\varepsilon_{0}>0$, $\rho$ and fixed faulty fraction
$\kf/\n$, the right-hand side of~\eqref{eq:pmiss} is strictly decreasing
in~$\n$.
\end{Corollary}

Corollary~\ref{cor:pmiss-monotone} is the safety half of the trade-off: it is
the reason a larger validator set is desirable, and it opposes the throughput
law of Section~\ref{sec:bifactor}, which is strictly decreasing in~$\n$ as
well --- but in the wrong direction for capacity. Section~\ref{sec:sizing}
exploits exactly this opposition.

\begin{Corollary}[Minimum validator count for a safety target]
\label{cor:nmin}
Requiring $\pmiss(\n,\w)\le\epss$ is equivalent to
\begin{equation}
\n\ \ge\ \nmin
\eqdef\frac{\bigl(1+(\kf-1)\rho\bigr)\,\ln(1/\epss)}
            {2\,\w\,\varepsilon_{0}^{2}} .
\label{eq:nmin}
\end{equation}
\end{Corollary}
